\subsection{Our Motivation}
Graphs are fundamental data structures that provide powerful tools for modeling relationships between various entities. However, there are computational are limitations on the scale of a graph, both physically and practically. Physically, a graph with many individual and dense clusters that are connected by very sparse regions may require large amounts of computational power. Practically, data is not always stored in one central graph and is often segmented into various subgraphs. For example, individual airlines may have their own, often dense, graph of their routes, however, they may not have the routes of other airlines.  Expanding on the previous example, while this may not be a problem for traveling most routes, say there doesn’t exist a route between city A and city B in any singular airline network. How can we find if it is possible to use air travel from city A to city B?


We can check to see if there exists an airline network that contains city A and an airline network that contains city B that share a common airport, or node on the graph. This practice is already being employed, many flights span across different airlines. What we seek to study is the amount of communication or information that each airline must share with one another to solve various problems. More generally, we seek to study the communication complexity of graph ownership problems.
